\section{Comparison with other balanced search trees}
\label{sec:comparison}

The AVL tree was the first structure to keep a search tree balanced, and it is
not the only one. The alternatives differ in how strictly they constrain the
shape of the tree, and that single choice decides the rest: a stricter invariant
buys shorter trees and faster searches, and it charges for them in rebalancing
work on every update. \Cref{tab:comparison} sets out where each structure lands.

\begin{table}[t]
  \centering
  \caption{Balanced search trees compared. Heights are worst-case bounds for $n$
  keys. The AVL tree keeps the shortest tree and pays for it during updates; the
  red-black tree accepts a tree up to twice as tall in exchange for cheaper
  rebalancing.}
  \label{tab:comparison}
  \begin{tabular}{@{}lllll@{}}
    \toprule
    Structure & Height bound & Rotations per insert & Rotations per delete
      & Suits \\
    \midrule
    AVL         & $1.44\lgg n$        & $\le 2$    & $O(\log n)$
      & search-heavy work \\
    Red-black   & $2\lgg(n+1)$        & $\le 2$    & $\le 3$
      & update-heavy work \\
    Splay       & amortised $O(\log n)$ & $O(\log n)$ & $O(\log n)$
      & skewed access \\
    B-tree      & $O(\log_B n)$       & none       & none
      & block storage \\
    \bottomrule
  \end{tabular}
\end{table}

\subsection{Red-black trees}

A red-black tree colours its nodes and requires that every root-to-leaf path
contain the same number of black nodes, with no two red nodes adjacent. The
condition allows one path to be twice as long as another~\cite{clrs}, so the height bound is
$2\lgg(n+1)$ against the $1.44\lgg n$ of \Cref{thm:heightbound}. Searches
therefore cost more in the worst case.

The compensation appears during deletion. \Cref{thm:deletebound} allows an AVL
deletion to rotate at every level, and the cascade of \Cref{fig:cascade} shows
the bound is attained. A red-black deletion needs at most three rotations
regardless of the size of the tree, because the looser invariant can absorb an
imbalance by recolouring rather than by restructuring, and recolouring may
propagate upward without moving any pointers.

Which structure wins depends on the ratio of searches to updates, and the
measured difference is smaller than the bounds suggest. Pfaff compared the two
structures across seventy-nine workloads and found the ratio of AVL running time
to red-black running time ranging from $0.677$ to $1.077$, with a median of about
$0.947$ and a geometric mean of about $0.910$~\cite{pfaff2004}.
\Cref{fig:pfaffrange} draws that spread. The AVL tree was faster in the majority
of those measurements, and neither structure dominated: the interval straddles
parity, so the workload decided the winner in both directions.

\begin{figure}[t]
  \centering
  % Pfaff's AVL-to-red-black running-time ratios, drawn from the figures reported
% in the paper rather than reproduced from it.
%
% The source paper is under publisher copyright, so its own figures cannot be
% reused here. The measurements themselves are facts and are citable, so this
% is an independent rendering of the reported range, median and geometric mean,
% credited in the caption. Do not replace it with a scan.
%
% Axis: ratio of AVL time to red-black time. Below 1 the AVL tree is faster.
\begin{tikzpicture}[x = 1mm, y = 1mm]

  \def\axl{0.62}   % left end of the axis, in ratio units
  \def\axr{1.16}
  % ratio -> millimetres: 200mm per unit ratio, so 0.54 units spans 108mm
  \newcommand{\rx}[1]{(#1 - \axl) * 200}

  % ---------------------------------------------------------------- the band
  \fill[rwash]
    ({\rx{0.677}}, -3.2) rectangle ({\rx{1.077}}, 3.2);
  \draw[draw = rfaint, line width = 0.4pt]
    ({\rx{0.677}}, -3.2) rectangle ({\rx{1.077}}, 3.2);

  % ------------------------------------------------------------- parity line
  \draw[draw = rkink, line width = 0.7pt, dashed]
    ({\rx{1.0}}, -12) -- ({\rx{1.0}}, 9);
  \node[tlabel, above] at ({\rx{1.0}}, 9) {equal speed};

  % ----------------------------------------------------------- the two marks
  \draw[draw = rkink, line width = 1.1pt]
    ({\rx{0.910}}, -3.2) -- ({\rx{0.910}}, 3.2);
  \node[tlabel, below] at ({\rx{0.910}}, -4.6) {geometric mean $0.910$};

  \draw[draw = rmid, line width = 1.1pt]
    ({\rx{0.947}}, -3.2) -- ({\rx{0.947}}, 3.2);
  \node[tlabel, above] at ({\rx{0.947}}, 4.2) {median $0.947$};

  % ----------------------------------------------------------- range end tags
  % the offset has to sit inside the expression group, not after it
  \node[tlabel, left]  at ({\rx{0.677} - 1.5}, 0) {$0.677$};
  \node[tlabel, right] at ({\rx{1.077} + 1.5}, 0) {$1.077$};

  % ------------------------------------------------------------------- axis
  % Sits at -12 rather than just under the band, so the geometric-mean label
  % at -4.6 has clear space between the band and the tick row.
  \draw[draw = rmid, line width = 0.5pt] ({\rx{\axl}}, -12) -- ({\rx{\axr}}, -12);
  \foreach \t in {0.7, 0.8, 0.9, 1.0, 1.1} {
    \draw[draw = rmid, line width = 0.5pt] ({\rx{\t}}, -12) -- ({\rx{\t}}, -13.6);
    \node[tlabel, below] at ({\rx{\t}}, -13.8) {\t};
  }
  \node[tannot, below] at ({\rx{0.89}}, -19)
    {ratio of AVL running time to red-black running time};

  % ------------------------------------------------------ which side is which
  \draw[-{Stealth[length=2.2mm]}, draw = rmid]
    ({\rx{0.99}}, 14) -- ({\rx{0.90}}, 14);
  \node[tlabel, left] at ({\rx{0.90}}, 14) {AVL faster};
  \draw[-{Stealth[length=2.2mm]}, draw = rmid]
    ({\rx{1.01}}, 14) -- ({\rx{1.10}}, 14);
  \node[tlabel, right] at ({\rx{1.10}}, 14) {red-black faster};

\end{tikzpicture}

  \caption{Relative running time of AVL against red-black trees across
  seventy-nine workloads. Drawn from the range, median and geometric mean
  reported by Pfaff~\cite{pfaff2004}; the rendering is ours.}
  \label{fig:pfaffrange}
\end{figure}

Brown reports more recent timings against three red-black variants~\cite{brown2025}.
The bottom-up red-black tree inserts and deletes randomly ordered keys faster
than the AVL tree, while the AVL tree is faster at inserting consecutively
ordered keys and slower at deleting them. The left-leaning variant was slower
than the other three on every workload tested. The practical conclusion is that
the choice between AVL and red-black trees is decided by the workload and the
implementation rather than by the height bounds.

\subsection{Splay trees}

A splay tree~\cite{sleator1985} abandons the idea of an invariant. After every access it rotates the
accessed node to the root, so the structure carries no balance information and
gives no worst-case guarantee for a single operation: one access can cost
$\Theta(n)$. What it guarantees is amortised, namely that any sequence of $m$
operations costs $O(m \log n)$.

The design pays off when accesses are skewed. A splay tree keeps recently used
keys near the root, so a workload that repeatedly touches a small subset
outperforms any structure that treats all keys alike, including the AVL tree.
The design fails when a guarantee is needed for each individual operation, as in
a real-time system, where an amortised bound is not usable.

\subsection{B-trees}

A B-tree~\cite{bayer1972} stores many keys in each node and keeps all leaves at the same depth.
Its height is $O(\log_B n)$ for a branching factor $B$, which for a node sized to
a disk block reduces the depth of a large index to a handful of levels.

The comparison with the AVL tree is not really a comparison of balancing
strategies. A B-tree is designed for storage where the cost is dominated by the
number of blocks fetched rather than the number of keys compared. An AVL tree of
$10^{6}$ keys has height $23$ by \Cref{fig:plotheight}, so a search touches $24$
nodes; if each node is a separate disk block, that is $24$ block reads. A B-tree
with a branching factor of $100$ holds the same keys in three levels. The AVL
tree compares fewer keys and reads far more blocks, which is why database indexes
and file systems use B-trees while in-memory structures do not.
